\section{Complete Component and Support Inventory}

This inventory groups source files by their role instead of reproducing raw
captures or generated event arrays as prose. Module paths in the first table
are relative to \path{src/lm_idnet/}. Package \path{__init__.py} files establish
import namespaces; they are not additional detector algorithms.

\begin{longtable}{>{\raggedright\arraybackslash}p{0.43\textwidth}>{\raggedright\arraybackslash}p{0.49\textwidth}}
\caption{Installed Python implementation map.}\label{tab:source-inventory}\\
\toprule\textbf{File} & \textbf{Responsibility} \\
\midrule\endfirsthead
\toprule\textbf{File} & \textbf{Responsibility} \\
\midrule\endhead
\path{cli.py} & Command parser, dispatch, UTC/rotating logs, structured errors, and stable exit codes. \\
\path{config.py} & Strict shared detector, research, and partition-free edge configuration. \\
\path{artifacts.py} & Typed loading, fingerprints, atomic snapshots, locks, history, restoration, and rollback. \\
\path{model_fingerprint.py} & Deterministic inference-parameter and input-contract identity. \\
\path{dataset_fingerprint.py} & Streaming source-byte and data-shaping-policy identity. \\
\path{partitioning.py} & Purpose-scoped fit, calibration, development, and reserved final selections. \\
\path{randomness.py} & Independent named simulation and model-fitting generators. \\
\path{runtime.py} & Bounded JSON-line input, chronology, synchronous events, and signal lifecycle. \\
\path{smoke.py} & Public synthetic infrastructure experiment and subprocess reproducibility. \\
\path{algorithms/dirichlet.py} & Minka fixed-point fitting and fit diagnostics. \\
\path{algorithms/dirichlet_multinomial.py} & Multinomial coefficient, complete window probability, and raw/normalised scores. \\
\path{algorithms/log_likelihood.py} & Validated SciPy and native LM parameter-dependent matrix kernels. \\
\path{algorithms/lm_native.py} & ctypes loading, coefficient initialisation, contiguous inputs, native lock, and precision errors. \\
\path{algorithms/estimator_factory.py} & One configured construction path for the existing estimator. \\
\path{models/modeling_stage.py} & Fit-capture validation, training matrix, initial estimation, and model provenance. \\
\path{models/calibration_stage.py} & Benign calibration windows, score quantile/IQR, and fixed packet-count limit. \\
\path{models/scoring_stage.py} & Events, shared one-window detector, both update policies, and batch reports. \\
\path{models/schemas.py} & Models, thresholds, events, accepted snapshots, pointer, and bounded history schemas. \\
\path{models/forecasting_stage.py} & Synthetic forecast scaffold; not held-out future-traffic validation. \\
\path{processing/categories.py} & Canonical taxonomy and exclusive packet classification. \\
\path{processing/inventory.py} & Configured sources, checksums, complete classic-PCAP counts, and duplicate explanations. \\
\path{processing/pcap_validation.py} & Structural EOF validation, bounded warning examples, and corruption reports. \\
\path{processing/pcap_processor.py} & Raw-frame MAC/ARP filtering, supported Ethernet decoding, and timestamps. \\
\path{processing/packet_transformer.py} & UTC time series, aligned windows, silence, discontinuities, and count matrices. \\
\path{processing/processing_manager.py} & Configured capture orchestration and processed metadata. \\
\path{processing/schemas.py} & Count/type/state validation and shared category/duration contract. \\
\path{processing/storage.py} & Validated canonical processed JSON input/output. \\
\path{processing/statistics.py} & Per-capture descriptive counts and daily dispersion fits. \\
\path{processing/timestamps.py} & Explicit timezone-aware UTC normalisation. \\
\path{processing/window_policy.py} & Window duration and common half-open alignment policy. \\
\path{evaluation/evaluation_stage.py} & Development event/label alignment, confusion matrix, FPR, and summary metrics. \\
\path{evaluation/benchmark_stage.py} & Isolated LM/SciPy fit, matrix, score, and decision comparisons. \\
\path{evaluation/metrics.py} & TVD, KL, MSE, Brier, and Jensen--Shannon helpers. \\
\path{evaluation/schemas.py} & Experiment-manifest provenance schema. \\
\path{exceptions.py} & Domain failure classes, committed-state distinction, and CLI exit mapping. \\
\bottomrule
\end{longtable}

\subsection{Native and supplied scientific files}

The production native directory contains \path{LM_time_lib.c}, horizontal-shift
and initialisation C/header files, the constants header, coefficient tables,
\path{LICENSE.}, and the shared object. The image builds its own shared object;
a checked-in binary is not trusted as architecture-independent evidence.
Only the required runtime library and log-Gamma coefficient files need to be
installed. Related digamma/trigamma tables remain supplied material, not a
claim that the Python alpha update has been replaced by a new native solver.

The preserved capsule under \path{lm-algorithm/} contains the original and
instrumented Python comparison drivers, executors, source and compiled C
library, coefficient tables, make/run helpers, row-scrambling helper, nine
cross-domain example CSV datasets, and fitting/test CSV inputs. Its environment
Docker/post-install files, metadata YAML, and \path{REPRODUCING.md} document the
original Code Ocean capsule. These are third-party reproduction assets. The
production package and the local paper-shaped experiment have their own
contracts, rather than silently importing every capsule helper.

\subsection{Project and dataset scripts}

\begin{longtable}{>{\raggedright\arraybackslash}p{0.43\textwidth}>{\raggedright\arraybackslash}p{0.49\textwidth}}
\caption{Executable research, deployment, and preparation paths.}\label{tab:script-inventory}\\
\toprule\textbf{File} & \textbf{Responsibility} \\
\midrule\endfirsthead
\toprule\textbf{File} & \textbf{Responsibility} \\
\midrule\endhead
\path{scripts/run.sh} & Build/run convenience for research commands with persistent data/result mounts. \\
\path{scripts/run-edge.sh} & Prebuilt offline constrained streaming/rollback launcher. \\
\path{scripts/docker-entrypoint.sh} & Direct installed-command execution; no startup compilation. \\
\path{scripts/benchmark_edge.py} & Finite synthetic replay instrumenting the real runtime, fitting, resources, and hashes. \\
\path{scripts/lm_professor_experiment.py} & Paired paper-shaped aggregate likelihood and forecast comparison. \\
\path{scripts/evaluate_window_sizes.py} & Matched labelled duration study with fresh models, time-scaled adaptation, annotation coverage, delay, and saved decisions. \\
\path{scripts/latex.sh} & Build technical LaTeX reports through the existing Compose service. \\
\path{scripts/run_adaptation_experiments.sh} & Select six conservative or six legacy configurations, adapt, then evaluate. \\
\path{unsw_iot_attack_pcaps/download_pcaps.sh} & Resume selected device/date source downloads. \\
\path{unsw_iot_attack_pcaps/download_annotations.sh} & Retrieve the separate official device annotations. \\
\path{unsw_iot_attack_pcaps/prepare_dataset.sh} & Build flat symbolic-link source views without copying capture bytes. \\
\path{unsw_iot_attack_pcaps/prepare_labels.py} & Align official attack intervals to actual capture windows and record provenance. \\
\bottomrule
\end{longtable}

\subsection{Tests and fixtures}

The test modules are grouped below. Temporary files and subprocesses are
isolated by shared fixtures in \path{tests/conftest.py}; there is no need to
modify real detector artefacts for a fault test.

\paragraph{Configuration and CLI} \path{test_config.py}, \path{test_cli.py}, \path{test_failure_semantics.py}, \path{test_testing_architecture.py}.

\paragraph{Data identity and input processing} \path{test_capture_inventory.py}, \path{test_pcap_validation.py}, \path{test_dataset_fingerprint.py}, \path{test_partitioning.py}, \path{test_timestamps.py}, \path{test_processing.py}.

\paragraph{Numerical and artefact contracts} \path{test_algorithms.py}, \path{test_dirichlet_multinomial.py}, \path{test_likelihood.py}, \path{test_lm_native_matrix.py}, \path{test_estimator_factory.py}, \path{test_artifact_schemas.py}.

\paragraph{Training, calibration, scoring, metrics} \path{test_modeling_stage.py}, \path{test_calibration_stage.py}, \path{test_scoring_stage.py}, \path{test_evaluation_stage.py}, \path{test_benchmark_stage.py}.

\paragraph{State and deployment} \path{test_detector_state.py}, \path{test_runtime.py}, \path{test_edge_deployment.py}, \path{test_adaptation_runner.py}.

\paragraph{Dataset preparation and smoke} \path{test_unsw_label_preparation.py}, \path{test_unsw_scripts.py}, \path{test_smoke_fixture.py}.

\paragraph{Research comparisons} \path{test_detection_experiments.py}, \path{test_window_size_experiments.py}.


Fixture files include the tiny packet CSV, synthetic smoke JSON and its licence
README, five-window edge JSON lines, quarantined truncated-PCAP hex, and a
normalised configuration snapshot. They exercise exact contracts rather than
redistributing the full mixed-device packet archive.

\subsection{Configuration, data, artefact, and output inventory}

There are three base device configurations and one edge configuration. Six
conservative and six explicit legacy experiment files cover both labelled
devices and their threshold/refit variants. Current experiment paths separate
policies. Older adaptive output directories remain historical records rather
than being overwritten under new policy names.

At the pre-report snapshot, tracked data comprise 64 D-Link raw captures,
12 D-Link processed captures, 17 Chromecast processed captures, and twelve
Samsung processed captures. The label directories contain 18 Chromecast and
13 Samsung files, including their manifests. The configured processed counts
and role totals appear in the dataset section. Raw UNSW staging is ignored;
the prepared local symbolic-link views are not a second source archive.

Tracked reference artefacts include three initial model/threshold/event sets,
six historical adaptive event sets, six explicit legacy event sets, and the
prepared edge pointer/snapshot. Mutable local edge state is ignored. Retained
reports include capture checks/statistics, fingerprints, smoke evidence,
production and paper-shaped numerical benchmarks, static metrics, historical
and explicit legacy adaptations/evaluations, legacy reproduction evidence,
and the four edge measurement/provenance files. Not every possible configured
output exists: the six schedule-specific conservative comparison directories
remain incomplete. The central detection records retain cross-device static,
conservative, and legacy replays and later audits.

\subsection{Root build and collaboration files}

\path{Dockerfile}, \path{pyproject.toml}, \path{main.py}, and \path{README.md}
define the installed package and its user entry points. The two Compose files
separate edge operation from LaTeX tooling. Git and Docker ignore rules,
attributes, the commit-instruction location, and local installation/environment
files support development. Local environment contents are deliberately excluded
from the report. No new automated CI or hidden external service is implied by
those supporting filenames.

\subsection{Evidence revision map}

The source revisions used for the report's distinct evidence boundaries are:
\begin{center}
\small
\begin{tabular}{lll}
\toprule
\textbf{Source} & \textbf{Revision} & \textbf{Role} \\
\midrule
Edge detector snapshot & \texttt{db2a736} & Recorded simulation boundary \\
Thesis workspace snapshot & \texttt{0d66bf8} & Writing structure \\
\texttt{feature/adapt} & \texttt{cb8dd66} & Historical adaptive experiments \\
\texttt{feature/lm-loglikelihood} & \texttt{96fb241} & Numerical/tooling history \\
\texttt{feature/labeled-dataset} & \texttt{210f922} & Labelled-data/history material \\
\texttt{feature/temporal-model} & \texttt{b0aa35d} & Separate chronology/fold research \\
Earlier feature experiments & \texttt{4279d84} & Historical detector research \\
Retained detection/scoring & \texttt{99bc189} & Joint count scoring and update references \\
\bottomrule
\end{tabular}
\end{center}

\subsection{Thesis source inventory}

The writing workspace contains the main document, metadata, preamble,
PDF metadata sidecar, executable build script, reference README, downloaded
upstream source, seal asset, eight front-matter files, seven chapter files,
bibliography source/output wrapper, and reproducibility appendix. The previous
section explains what each group does. Generated PDFs and auxiliary files
belong under ignored \path{out/}, not inside chapter source directories.
